\documentclass[review,12pt,number]{elsarticle}
\usepackage{amsmath,amssymb,booktabs,graphicx}
\usepackage[colorlinks=true,linkcolor=blue,citecolor=blue,urlcolor=blue]{hyperref}

\begin{document}

\begin{frontmatter}

\title{The Repair Saturates the Gate: Context Noise as a Degeneracy
Diagnostic in Active-Dendrite Continual Learning}

\author[hnucm]{Ranxi Zhou}
\affiliation[hnucm]{organization={Hunan University of Chinese Medicine},
            city={Changsha}, postcode={410208}, country={China}}

\begin{abstract}
Active-dendrite continual-learning systems condition each neuron's
plasticity and routing on a task-context signal. A companion audit
(\cite{zhou2026delpS1}, S1) found that such networks never develop
segment--task specialization and that a constant bias $b$ on the
own-segment dendritic input causally repairs the learning term
($\Delta\ell$ up to $+4.6$, $|t| \ge 8$). Every result in that line --
including the repair itself -- was obtained under oracle context. We ask
the question the line left open: \emph{when the context signal is wrong,
what survives?} Under the worst-case evaluation protocol (each sample's
context one-hot flipped to a uniformly random other task with probability
$\rho$), the answer is: everything, at every noise level tested.
Final accuracy moves by at most $0.25$ points at $\rho \le 0.4$ for the
base system and by at most $0.01$ points for the repaired one; the
repair's paired gain ($+4.80$, $t{=}17.0$, on MNIST) is retained at
100--105\% under noise on MNIST and 100--124\% on Fashion-MNIST. Zero
noise sensitivity is not robustness --- it is a symptom. Mechanism
probes show why: the trained dendritic gate is \emph{context-degenerate}.
Its value under the correct context and under every wrong context is
identical to four decimals ($0.565$); after the repair ($b{=}10$)
context swaps change at most $0.02\%$ of predictions on average (worst
pair $0.15\%$); and the dendritic weight columns stay mutually
orthogonal (cosine $\approx 0.003$) throughout training. Two degenerate
controls then locate what the repair actually repairs. A
\emph{constant-gate} control that pins the gate at its measured baseline
value ($0.565$) reproduces the base network, not the repair --- merely
holding the gate constant buys nothing, because the base system already
runs at that constant. A \emph{no-gate} control ($h = \mathrm{kwta}(t)$,
same trunk, same anchor chain) recovers most of the repair's benefit on
MNIST ($67.50$ vs.\ $68.67$, $-1.16$, $t{=}-2.7$) and numerically
exceeds it on Fashion-MNIST ($44.27$ vs.\ $42.89$, n.s.): the repair
functions as a saturating soft deletion of the gate ($\sigma(10{+}u)
\approx 1$), and what it buys is the removal of a constant
${\approx}0.57$ damping factor --- not restored routing. We conclude
that ``context-gated protection'' in this architecture family runs, at
these training budgets, on a gate that does not gate by context;
noise-robustness certificates for such systems should therefore be read
as \emph{degeneracy diagnostics}, and repairs that survive adversarial
context corruption deserve particular scrutiny.
\end{abstract}

\begin{keyword}
continual learning \sep active dendrites \sep context gating \sep
robustness evaluation \sep measurement validity \sep catastrophic
forgetting
\end{keyword}

\end{frontmatter}

\section{Introduction}
\label{sec:intro}

Active-dendrite architectures give each neuron a bank of dendritic
segments and let a task-context signal select which segment gates the
neuron's output
\cite{iyer2022,boahen2022,poirazi2020}. In continual learning the
arrangement is attractive for a specific reason: if each task's context
activates its own segments, protection can be localized to exactly the
parameters that task used, and interference across tasks should shrink.
A sequence of audits from the same laboratory has stress-tested that
promise. The original DeLP audit\footnote{``Where do the gains come
from?'', \cite{zhou2026audit}; the precursor study.} corrected its own
baselines and found most of the reported advantage was calibration, not
mechanism. The companion dose--response audit (S1,
\cite{zhou2026delpS1}) found the specialization premise itself fails:
segment--task routing never specializes, the standard routing instrument
is contaminated by dead neurons, and yet a one-parameter intervention ---
a constant bias $b$ on the own-segment dendritic input --- causally
improves the learning term in every cell of a
$\{$dataset$\}\times\{$depth$\}\times\{$granularity$\}$ grid. A transfer
study (S3, \cite{zhou2026repair}) then showed the repair does not survive
stronger backbones. One assumption has been shared by every experiment in
this line, and by most of the active-dendrite continual-learning
literature: \emph{the context signal at evaluation is the correct one}.

This paper removes that assumption and reports what happens. The probe is
deliberately adversarial: at evaluation time, each sample's context
one-hot is flipped to a uniformly random \emph{other} task with
probability $\rho \in \{0.1, 0.2, 0.4\}$. Training is untouched. Three
questions, preregistered before any data was collected
(Sec.~\ref{sec:design}): (Q1)~does the base gated system degrade?
(Q2)~does the S1 repair --- whose premise is precisely that the gate
should route by context --- survive context corruption? (Q3)~does a soft
mixture mitigation help?

The empirical answers are monotone: nothing degrades, at any $\rho$, on
either dataset, with or without the repair
(Sec.~\ref{sec:results}). A robustness result this clean is not a
reason to celebrate; it is a reason to dig. The rest of the paper is the
dig. Mechanism probes (Sec.~\ref{sec:mech}) show the trained gate is
\textbf{context-degenerate}: its value is identical to four decimals
under every context, and after the repair the network's predictions are
invariant to context swaps up to $0.02\%$ of samples. Two degenerate
controls (Sec.~\ref{sec:ablation}) then locate what the repair actually
does: pinning the gate at its own baseline constant reproduces the base
network --- flattening buys nothing, because the base system already
runs at that constant --- while deleting the gate outright recovers most
of the repair's benefit on MNIST and numerically all of it on
Fashion-MNIST. The repair works by \emph{saturating the gate to a
constant gain of one}: a soft deletion of the system's dependence on the
component it was supposed to be fixing.

The contribution is threefold:
\begin{enumerate}
\item \textbf{A negative-space result with a warning attached.}
Adversarial context corruption is free at every noise level --- because
the mechanism it corrupts is not load-bearing at these training budgets.
Any ``robust to context noise'' certificate earned by such a system is
evidence of degeneracy, not of robustness.
\item \textbf{A mechanism re-qualification of the S1 repair.} The
own-segment bias does not restore segment--task routing (which S1 showed
never existed); it saturates the gate into a context-free gain of one ---
a soft deletion of the dendritic damping. A constant-gain control at the
measured baseline gate value reproduces the base system, and a
gate-deletion control reproduces most (MNIST) or more (Fashion-MNIST) of
the repair: what helps is \emph{removing} the constant ${\approx}0.57$
damping, not shaping it.
\item \textbf{A reusable three-instrument protocol}: worst-case noise
sweep $\to$ flip-rate/gate-response probe $\to$ degenerate-control
ablation with two variants (constant-gate and gate-removal). Each
instrument is cheap; together they convert a suspicious null into a
mechanistic verdict. We preregistered the first two, and the controls
were forced on us by their result.
\end{enumerate}

\section{Background and setting}
\label{sec:background}

\paragraph{Architecture.} The reference network is the one shared by the
S1/S3 line \cite{zhou2026delpS1,zhou2026repair}: a one-hidden-layer
active-dendrite network (ADNet) with $n_h{=}256$ hidden units, one
dendritic segment bank per task ($K{=}T$), $k$-winners-take-all sparse
output ($k{=}64$). Given input $x$ and context one-hot $c$:
\begin{equation}
t = W_{\mathrm{ff}} x, \qquad
u_k = w_k^{\top} c \;\; (k\text{-th segment row}), \qquad
h = \mathrm{kwta}\!\left(t \odot \sigma(d)\right), \quad
d = \max_k u_k,
\label{eq:adnet}
\end{equation}
i.e.\ the winning segment's scalar $d$ gates every hidden unit
multiplicatively. The repair of S1 adds a constant $b$ to the own
segment's $u_{\mathrm{own}}$ before the max, everywhere (train, eval,
usage) \cite{zhou2026delpS1}.

\paragraph{Protection.} Usage-localized elastic consolidation (DeLP
\cite{zhou2026audit}): after each task, a Fisher anchor is stored and
weighted by measured per-neuron usage; the penalty pulls shared
parameters toward anchored values with strength $\lambda{=}1000$
(matching S1's dose--response setting). Segment-specific gradients are
protected so each task trains its own dendritic block.

\paragraph{Protocol.} Permuted-MNIST and Permuted-Fashion-MNIST, $T{=}10$
tasks, 10 seeds, paired across all arms; permutations drawn once per seed
and shared. Metrics: final average accuracy $\bar{a}$ over all tasks
after the last task; paired $t$-statistics over seeds.

\paragraph{Execution correction (2026-10-27).} An independent
data-consistency audit of this line, run after the first version of these
experiments had completed, found two defects in the harness and the
experiments were re-run in full. (i)~The elastic penalty was applied
without its $\lambda$ factor, so every run had effectively executed with
$\lambda \approx 1$ while the record claimed $\lambda{=}1000$; the
penalty now carries $\lambda$ as the reference implementation does.
(ii)~The degenerate control had a hard-coded hidden width of $2000$ while
the dendritic arms use $256$, so the two arms of the ablation differed by
roughly eight-fold in parameter count; both now use the same width. The
numbers reported here are all from the corrected re-run; the first
version's records are retained in the artifact under an explicit
deprecation notice. The correction moves the absolute accuracies
substantially --- a genuinely enforced penalty forgets more, and the
base network drops ${\sim}7.6$ points on MNIST --- and it changes one
qualitative verdict: the first version's MNIST gate-removal
\emph{equivalence} was an artifact of the width mismatch, and under
matched widths the equivalence becomes approximate rather than exact
(Sec.~\ref{sec:ablation}). The robustness null and the mechanism probes
are unaffected.

\section{Preregistered design}
\label{sec:design}

The design was locked on 2026-10-25, before any data was written to disk
(preregistration document included in the artifact). Noise is injected
\emph{at evaluation only}; training sees oracle contexts. The noise model
is worst-case in a specific sense: with probability $\rho$ per sample the
context one-hot is replaced by a uniformly random \emph{other} task ---
never the correct one --- so corrupted samples exercise the gate under a
plausible wrong context rather than a neutral one.

\begin{itemize}
\item \textbf{E1 (gating robustness).} $b{=}0$, $\rho \in
\{0, 0.1, 0.2, 0.4\}$, both datasets. Survival: $\rho{=}0.1$ costs
$\le 2$ points and $\rho{=}0.2$ costs $\le 5$, \emph{and} the
per-seed direction agrees across both datasets in $\ge 8/10$ seeds;
death: $\rho{=}0.1$ costs
$>5$ or $\rho{=}0.2$ collapses ($\bar a < 50$ or cost $>10$).
\item \textbf{E2 (repair survival).} Same grid crossed with
$b \in \{0, 10\}$, paired per $(\rho, \text{dataset}, \text{seed})$.
Survival: at $\rho{=}0.1$ the repair's paired $\Delta\bar a > 0$ and
retains $\ge 50\%$ of its $\rho{=}0$ gain; death: $\Delta\bar a \le 0$.
\item \textbf{E3 (soft mixture, exploratory).} Context replaced by
$(1-\rho)\,\mathrm{onehot} + \rho\,\mathrm{uniform}$; direction only,
no survival claim (multiple-comparison hygiene).
\end{itemize}

Two deviations from the pre-registered wording, both forced and both
reported rather than absorbed. First, E2's criterion was written for
$\Delta\ell$; the noise harness records only the final-phase
per-task accuracies and no $\ell$ field, so $\Delta\ell$ is not
computable from the collected data and the criterion is evaluated on
$\Delta\bar a$ instead --- the same "net benefit of the repair"
quantity, in the one form the runs support. Second, E1's third
condition (per-seed directional agreement) was not met in strict
form: the per-seed $\rho$-effect has mixed signs, though every
per-seed shift stays below $0.1$ points. We report that as a
condition we did not satisfy rather than dropping it, and the two
amplitude conditions pass with orders of magnitude to spare.

The noise stream is fixed per seed (a dedicated generator), so paired
comparisons share corruption patterns exactly. The protocol trains one
model per $(\text{dataset}, b, \text{seed})$ and evaluates it at all
four $\rho$ levels: \textbf{60 training runs} (40 for E1+E2, 20 for E3)
yielding 240 evaluation cells. (The pre-registration counted $\rho$ as
independent training runs --- 240 --- which over-counts by a factor of
four; the number of evaluated $(\rho, b, \text{dataset}, \text{seed})$
cells is unchanged.) A dated amendment documenting every deviation is
included in the artifact
(``prereg/S4\_ctxnoise\_预注册\_修订说明\_2026-10-27.md''); the original
pre-registration text is left untouched.

\section{Results: nothing degrades, at any noise level}
\label{sec:results}

\begin{table}[t]
\centering
\caption{Final average accuracy under worst-case context flips
($n{=}10$ seeds; noise applied to the final-phase evaluation only;
$\lambda{=}1000$ enforced, $n_h{=}256$ in all arms).}
\label{tab:main}
\begin{tabular}{llrrrr}
\toprule
dataset & arm & $\rho{=}0$ & $\rho{=}0.1$ & $\rho{=}0.2$ & $\rho{=}0.4$ \\
\midrule
MNIST    & base ($b{=}0$)    & 63.86 & 63.79 & 63.72 & 63.62 \\
MNIST    & repaired ($b{=}10$) & 68.67 & 68.67 & 68.67 & 68.67 \\
Fashion  & base ($b{=}0$)    & 42.20 & 42.15 & 42.09 & 42.04 \\
Fashion  & repaired ($b{=}10$) & 42.89 & 42.89 & 42.89 & 42.89 \\
\bottomrule
\end{tabular}
\end{table}

Table~\ref{tab:main} is the entire empirical content of E1 and E2.
Both preregistered survival criteria pass with orders of magnitude to
spare: the largest movement anywhere is $0.25$ points (MNIST base,
$\rho{=}0.4$), the largest on Fashion-MNIST is $0.16$, and the repaired
arm moves by at most $0.01$ points on MNIST and \emph{exactly} $0.000$
points on Fashion-MNIST --- where the death criteria allowed $10$ and
$5$. The repair's paired gain on MNIST is $+4.80$ ($t{=}17.0$) at
$\rho{=}0$ and $+5.05$ ($t{=}17.3$) at $\rho{=}0.4$; retention of the
gain is $100.0\% \to 105.1\%$ across the sweep --- under noise the
repair \emph{helps slightly more}, because its floor is flatter.
On Fashion-MNIST the gain is small, as S3 found at $T{=}10$
\cite{zhou2026repair}: $+0.69$ ($t{=}1.5$, n.s.) at $\rho{=}0$, rising
to $+0.86$ ($t{=}1.8$) at $\rho{=}0.4$ (retention $100\% \to 123.7\%$;
all n.s.). E1's third condition (per-seed directional agreement) holds
in strict form on MNIST --- $10/10$ seeds shift negatively at every
$\rho$ --- and fails on Fashion-MNIST ($8/10$, $9/10$, $7/10$ negative,
with every per-seed shift within $\pm 0.18$ points); we report it as
satisfied on one dataset and not on the other, as disclosed in
Sec.~\ref{sec:design}. E3's soft mixture tracks the oracle slightly
more closely than hard flips on the base arm (MNIST $\rho{=}0.4$: soft
$-0.09$ vs.\ hard $-0.25$; paired soft$-$hard $+0.16$, $t{=}3.8$) and
does nothing on the repaired arm; every E3 effect is ${\le} 0.26$
points.

A protocol-level artifact is worth excluding explicitly: the noise
\emph{is} reaching the network. Per-seed corruption patterns are fixed,
the context tensor is modified before the forward pass, and --- the
direct check --- swapping contexts \emph{does} change predictions in the
base network ($14.8\%$ of samples on average on MNIST, $30.1\%$ on
Fashion-MNIST, up to $51.7\%$ for particular swaps;
Sec.~\ref{sec:mech}). What is true is that prediction changes of this
kind average out across a $10$-way task mixture at the final-accuracy
level. The null is real, and it demands a mechanism.

\section{Mechanism: the gate does not gate by context}
\label{sec:mech}

Three probes, each applied to fully trained networks. The headline
numbers below are MNIST seed 0; the pattern replicates on Fashion-MNIST
seed 0 and MNIST seed 1 (all numbers in the artifact): at $b{=}0$,
flip rates $14.8$--$30.1\%$ mean / $23.8$--$51.7\%$ worst pair; at
$b{=}10$, flip rates $0.007$--$0.021\%$ mean / at most $0.15\%$ worst
pair; gate own/other identical to four decimals ($0.564$--$0.565$); and
mean column cosines in $[0.001, 0.003]$ in every case.

\paragraph{Prediction flip rate.} For each task $c$ and each alternative
context $c' \ne c$, measure the fraction of test samples whose argmax
prediction differs between $\mathrm{ctx}{=}c$ and $\mathrm{ctx}{=}c'$.
Base network: mean $14.8\%$ (MNIST) and $30.1\%$ (Fashion-MNIST), worst
pair up to $51.7\%$ --- the base output is genuinely, substantially
context-sensitive at the individual-prediction level, yet none of that
sensitivity survives to final accuracy (Sec.~\ref{sec:results}: at most
$0.25$ points). The context pathway does something per sample; what it
does nets out to nothing. Repaired network ($b{=}10$): mean
$0.007$--$0.021\%$, worst pair $0.15\%$ --- no longer \emph{exactly}
zero, but a rounding error: roughly one flipped prediction in five
thousand on average, a dozen samples for the worst single pair.

\paragraph{Gate response.} The winning segment's gate value
$\sigma(d)$ has mean $0.565$ under the correct context and $0.565$
under every alternative --- identical to four decimal places, in both
the base ($0.5651$) and repaired ($0.5652$) networks. The context
one-hot selects \emph{which} segment wins, but not \emph{how much} it
gates: the gate's value is a constant of the training run, not a
function of the context.

\paragraph{Column geometry.} The dendritic weights projected onto the
$T$ context columns remain mutually orthogonal after training: mean
pairwise cosine $0.003$ (per-model mean in $[0.001, 0.003]$; the most
negative single pair over the three models is $-0.056$), indistinguishable
from their random initialization. Training, with per-task gradient
protection, leaves the context channel of the gate bank essentially at
its prior.

Together: the architecture's context pathway is present, trainable in
principle, protected per task --- and unused as a \emph{modulator}. The
max over segments is won by whichever row the (orthogonal, near-random)
columns happen to favor, the gate value is the same constant
($\approx 0.565$) no matter which segment wins, and the repair adds a
constant so large ($b{=}10$ against $|u| \lesssim 1$) that the own
segment always wins with $\sigma(10 + u) \approx 1$: a context-free
gain of one. Equation~\ref{eq:adnet} degenerates, in both arms, to a
multiplication by a constant --- ${\approx}0.57$ in the base network,
${\approx}1.0$ in the repaired one. Which of those two constants
matters, and why, is what the next section measures.

\section{Ablation: what does the repair repair?}
\label{sec:ablation}

If the repaired network computes (approximately) $h = \mathrm{kwta}(t)$,
then a network with the dendritic bank deleted and the same trunk,
sparsifier, and anchor chain should reproduce most of the repair's
benefit --- and, conversely, a network whose gate is held at the base
system's measured constant ($0.565$) should reproduce the base system,
because the base gate already \emph{is} that constant
(Sec.~\ref{sec:mech}). We therefore run four arms under one protocol
($\rho{=}0$, $n{=}10$ paired seeds, all arms at $n_h{=}256$, all with
usage-weighted Fisher anchors at $\lambda{=}1000$): \textbf{base}
($b{=}0$), \textbf{repaired} ($b{=}10$), \textbf{no-gate}
($h = \mathrm{kwta}(W_{\mathrm{ff}} x)$; the bank deleted), and
\textbf{constant-gate} ($h = \mathrm{kwta}(0.565 \cdot
W_{\mathrm{ff}} x)$; the gate pinned at the base system's measured
mean). The base and repaired arms replicate the Sec.~\ref{sec:results}
runs to two decimals. The degenerate controls were added during the
post-audit re-run (Sec.~\ref{sec:background}); they were not
preregistered. The no-gate control exists in three versions in the
artifact --- an invalidated first attempt (an extra sigmoid on $t$ that
changed the activation structure), a width-mismatched second version
($n_h{=}2000$), and the matched version reported here --- all retained
under an explicit deprecation notice.

\begin{table}[t]
\centering
\caption{Four-arm control experiment ($n{=}10$ paired seeds,
$\rho{=}0$, matched width $n_h{=}256$ and matched usage-weighted
anchors). Bold marks the best arm per dataset.}
\label{tab:ablation}
\begin{tabular}{lrrrr}
\toprule
dataset & base ($b{=}0$) & const-gate ($0.565$) & repaired ($b{=}10$)
& no-gate \\
\midrule
MNIST   & 63.86 & 63.62 & \textbf{68.67} & 67.50 \\
Fashion & 42.20 & 42.70 & 42.89 & \textbf{44.27} \\
\bottomrule
\end{tabular}
\end{table}

\textbf{The constant-gate control reproduces the base system, not the
repair.} Pinning the gate at $0.565$ lands at $63.62$ on MNIST (base:
$63.86$; paired difference $-0.25$, $t{=}-0.7$, n.s.) and $42.70$ on
Fashion (base: $42.20$; $+0.49$, $t{=}0.7$, n.s.). Holding the gate
constant buys nothing --- which is exactly what the mechanism section
predicts, since the base network's gate already \emph{is} constant at
that value. The entire difference between the base and repaired
networks is the value of that constant: $0.57$ damps the trunk,
$1.0$ does not.

\textbf{Gate deletion recovers most of the repair --- and on
Fashion-MNIST, more than the repair.} The no-gate control reaches
$67.50$ on MNIST, recovering $+3.64$ of the repair's $+4.80$-point gain
($76\%$) and falling short of the repaired network by $-1.16$
($t{=}-2.7$): keeping the (saturated) multiplicative structure is worth
about one point, marginally significant at $n{=}10$. On Fashion-MNIST
the no-gate control is the \emph{best} arm: $+2.07$ over base
($t{=}2.5$) and $+1.37$ over the repaired network ($t{=}1.9$, n.s.) ---
on the harder dataset, full deletion beats the repair outright.

\textbf{Sensitivity to anchor weighting.} A second no-gate run with
plain (unweighted) Fisher anchors --- a stronger, mismatched chain ---
lands at $63.63$ on MNIST and $41.10$ on Fashion, i.e.\ at or below
base. Anchor weighting alone moves the no-gate arm by $3.9$ and $3.2$
points. The matched-anchor four-arm table above is therefore the
primary comparison, and the residual MNIST gap ($-1.16$) should be read
against that sensitivity. The first version of this experiment reported
a clean MNIST \emph{equivalence} between the no-gate control and the
repair; that verdict was an artifact of the width mismatch disclosed in
Sec.~\ref{sec:background} and is retracted.

\textbf{Verdict.} The repair is a saturating soft deletion: $b{=}10$
drives $\sigma \to 1$, which is the no-gate limit reached through the
gate rather than around it. What the repair buys on these benchmarks is
the removal of a constant ${\approx}0.57$ damping factor --- and on
Fashion-MNIST, where deletion is free to remove it \emph{exactly}, the
removal is worth slightly more than the repair itself. What does
\emph{not} follow is the stronger claim of the first version, that the
repair and gate removal are interchangeable on MNIST: with matched
widths and matched anchors they are close ($76\%$ of the gain
recovered) but distinguishable.

\section{Discussion}
\label{sec:discussion}

\paragraph{For the S1 story.} The dose--response audit established that
the own-segment bias repairs the learning term; the mechanism section of
that paper showed the network never pushes back at the parameter level.
This paper supplies the missing re-qualification: the repair does not
restore routing --- there was no routing to restore --- it \emph{saturates
the gate into a context-free gain of one}, i.e.\ a soft deletion of a
constant ${\approx}0.57$ damping factor. Read together, the chain is
now: no specialization (S1) $\to$ repair works (S1) $\to$ the repair is
a saturating gate deletion that recovers $76\%$ of the removal benefit
on MNIST and is beaten by full removal on Fashion-MNIST (this paper)
$\to$ the repair does not transfer to strong backbones (S3). What
survives as the ``mechanism'' of the repair is the removal of a constant
damping --- sparse capacity management at best, context-gated routing
not at all.

\paragraph{For evaluation practice.} Robustness certificates obtained
under context corruption are usually read as good news. In this family
they are a diagnostic: a system whose accuracy is invariant to
\emph{adversarial} context corruption is a system whose context pathway
is not load-bearing. The three-instrument protocol
(Sec.~\ref{sec:results}--\ref{sec:ablation}) costs a few GPU-hours and
converts that ambiguity into a verdict. We suggest it as a standing
companion to any context-gated continual-learning evaluation, in the
same spirit as the dead-neuron probe of S1.

\paragraph{Limitations.} (i) Noise is injected at evaluation only;
training-time context errors corrupt usage statistics and anchors
differently, and are the natural next experiment. (ii) The mechanism
probes are computed on fully trained networks at $T{=}10$; whether the
context columns stay orthogonal at larger $T$ or longer training is
untested. (iii) The degenerate controls were not preregistered and were
added after the audit; the anchor-weighting sensitivity of the no-gate
arm ($\pm 3.9$ points, Sec.~\ref{sec:ablation}) is larger than the
residual repair-vs-deletion gap it bounds. (iv) One architecture
family, two digit datasets --- the same scope limits as the companion
line.

\section{Conclusion}

Asked whether context-gated protection survives a wrong context signal,
the system answered: it never needed the signal. The repair that was
motivated by broken routing survives for the same reason --- it never
used the routing; it removes a constant damping factor that the
never-used gate had left in place, and where removal can be done
exactly, it does slightly better than the repair. The general lesson
travels beyond this architecture: \emph{when a robustness probe returns
exactly zero damage, measure the pathway before celebrating the armor
--- and when a repair works, check whether it works by fixing the
mechanism or by deleting it}.

\section*{Data and code availability}
All code, per-run records (including per-$\rho$ accuracy vectors), the
preregistration document, and both versions of the degenerate control
(including the invalidated first attempt) are in the artifact repository
\url{https://github.com/zhouranxi/delp-anchor-budget} (folder
\texttt{S4\_gate\_bypass/}, self-contained: paper sources, scripts, raw
records), alongside the companion papers S1--S3.

\section*{Declaration of competing interest}
The author declares no competing interests.

\bibliographystyle{elsarticle-num}
\bibliography{refs}

\end{document}
